%%%% CAPÍTULO — CONSIDERAÇÕES FINAIS
%%
\chapter{Considerações Finais}\label{cap:consideracoesfinais}

Este capítulo encerra a monografia: retoma os objetivos em face do que foi realizado, explicita limitações do escopo e indica caminhos de continuidade do E-Museu para trabalhos acadêmicos e técnicos posteriores. O processo de desenvolvimento está no Capítulo~\ref{cap:desenvolvimento}; a consolidação dos resultados, na Tabela~\ref{tab:objetivos-resultados} e nas demais seções do Capítulo~\ref{cap:resultados}.


\section{Síntese em relação aos objetivos}\label{sec:conclusao-objetivos}

O objetivo geral deste trabalho consistiu em refatorar e implementar novas funcionalidades no E-Museu, aplicando boas práticas de programação, padrões de projeto e testes automatizados, de modo a melhorar a manutenibilidade e a qualidade do código (Seção~\ref{subsec:objetivoGeral}).

Em relação aos objetivos específicos (Seção~\ref{subsec:objetivoEspecificos}), a Tabela~\ref{tab:objetivos-resultados} registra atendimento pleno a todas as metas: estrutura de versionamento no \textit{GitHub}, configuração de ambientes, refatoração do código legado, ampliação da suíte de testes automatizados de fluxo e integração (Seção~\ref{sec:dev-testes}), novas funcionalidades priorizadas no MoSCoW (internacionalização, múltiplas imagens por item, identificação com QR Code, Turnstile, tradução assistida no painel, entre outras, Seção~\ref{sec:dev-funcionalidades}) e aplicação de padrões e boas práticas descritos no referencial teórico.

Em síntese, o trabalho partiu das limitações da versão anterior \cite{gonzaga2024} e concentrou-se em estabelecer uma base técnica mais organizada para o projeto institucional, conforme os Capítulos~\ref{cap:desenvolvimento} e~\ref{cap:resultados}.


\section{Limitações}\label{sec:conclusao-limitacoes}

Este trabalho não esgotou todas as possibilidades de evolução do E-Museu. O escopo concentrou-se em refatoração, infraestrutura, testes e funcionalidades de acervo priorizadas no MoSCoW, com um único desenvolvedor e recursos de tempo e financeiros limitados (Seção~\ref{sec:dev-esforco}). A interface pública e o painel administrativo mantiveram, em grande parte, o visual da versão anterior (Seção~\ref{sec:resultados-interface}). O pacote de continuidade entregue à orientação (Seção~\ref{sec:dev-deploy-continuidade}) reduz o risco de perda de dados, mas não substitui rotinas automáticas de backup em produção.


\section{Trabalhos futuros e continuidade do E-Museu}\label{sec:conclusao-trabalhos-futuros}

O E-Museu é um projeto de longa duração, vinculado à \acrshort{utfpr} e a iniciativas de extensão relacionadas à preservação da história da informática. O presente trabalho deixou o sistema e o repositório em condições mais sustentáveis para quem der continuidade ao desenvolvimento, seja em novo Trabalho de Conclusão de Curso, seja em extensão ou na manutenção do sistema na \acrshort{utfpr}. O repositório institucional, a wiki de fluxo de trabalho, a documentação de implantação e o pacote de continuidade descritos no Capítulo~\ref{cap:desenvolvimento} servem de ponto de partida para não repetir a reorganização básica já feita neste ciclo.

A principal frente sugerida para trabalhos posteriores é a restilização da interface, pública e administrativa. O foco deste TCC foi a base técnica; as capturas em produção (Seção~\ref{sec:resultados-interface}) mostram que \textit{layout}, tipografia e navegação permanecem próximos aos de \citeonline{gonzaga2024}. Uma reforma visual planejada com a orientação poderia melhorar usabilidade, acessibilidade e percepção do museu virtual sem comprometer a arquitetura já entregue.

Em seguida, destacam-se evoluções operacionais e de gestão do acervo:

\begin{itemize}
    \item Backup automático do banco de dados e das imagens do acervo (por exemplo agendamento no servidor ou integração com o \textit{MinIO} e o \textit{MySQL}), complementando o pacote manual montado neste TCC e reduzindo dependência de exportações pontuais.
    \item Relatórios de acesso ao catálogo público e ao painel (visitas, rotas mais usadas, período), para apoiar decisões de extensão e de divulgação do museu.
    \item Registro de logs das alterações feitas por administradores (quem editou ou excluiu item, categoria ou colaborador, e quando), reforçando rastreabilidade da curadoria.
    \item Newsletter por e-mail, aproveitando a infraestrutura de mensagens transacionais já usada na verificação de colaboradores e na contribuição (Seção~\ref{sec:dev-func-email-redis}), para avisar interessados sobre novidades no acervo ou no projeto.
    \item Campo de link de vídeo de demonstração por item (URL externa), para enriquecer a ficha da peça sem armazenar o arquivo de vídeo no servidor.
\end{itemize}

Outras melhorias podem ampliar a confiabilidade e a qualidade percebida: mais testes automatizados (incluindo usabilidade ou acessibilidade), avaliação com usuários do acervo, monitoramento do ambiente \textit{Coolify} e redução de intervenções presenciais na máquina da \acrshort{utfpr}, quando a política institucional permitir.

Recomenda-se que futuros trabalhos formalizem requisitos e cronograma com a orientação, reutilizando o fluxo Gitflow documentado \cite{gitflowdoc}, de modo a acumular esforço em funcionalidades e em validação sobre a base técnica já consolidada.

Assim, encerra-se o Trabalho de Conclusão de Curso com o E-Museu em operação nos domínios institucionais, repositório organizado e objetivos específicos atendidos conforme documentado nesta monografia. O sistema segue disponível para visitantes e para a gestão do acervo; o texto registra o percurso técnico realizado para que orientação, docentes e quem assumir a manutenção possam dar continuidade ao museu virtual como espaço de preservação e divulgação da memória tecnológica na \acrshort{utfpr}.
